\documentclass[a4paper,12pt]{article}
\usepackage[dvipsnames]{xcolor}
\usepackage[skip=10pt plus1pt, indent=0pt]{parskip}
\usepackage{graphicx} % Required for inserting images
\usepackage{listings}
\usepackage{amsmath}
\usepackage{float}
\usepackage[italian]{babel}
\usepackage{geometry}
 \geometry{
 a4paper,
 left=25mm,
 right=25mm,
 top=30mm,
 bottom=30mm,
 }
 \usepackage{hyperref}
\hypersetup{
    colorlinks=true,
    linktoc=all,
    linkcolor=black,
    citecolor=black,
}

\title{Algoritmo parallelo di Smith-Waterman per allineamento di sequenze DNA}
\author{Christian Pozzoli 29678A}
\date{Ottobre 2025}

\begin{document}
\lstset{
language=C++,
basicstyle=\fontfamily{pcr}\footnotesize,
commentstyle=\color{OliveGreen}\ttfamily,
keywordstyle=\color{NavyBlue}\bfseries,
numberstyle=\tiny\color{Gray},
%ndkeywordstyle=\bfseries,
numbers=left
}
\maketitle
\begin{figure}[H]
    \centering
    \includegraphics[width=0.8\linewidth]{images/cover.jpg}
\end{figure}
\newpage
\renewcommand*\contentsname{Indice}
\tableofcontents
\newpage

\section{Introduzione}
\subsection{Allineamento di sequenze di DNA}
L'allineamento di sequenze di DNA è una delle operazioni fondamentali della bioinformatica e rappresenta uno strumento essenziale per lo studio comparativo dei genomi.
L'obiettivo principale dell'allineamento è confrontare due o più sequenze di basi di nucleotidi al fine di individuare delle regioni di somiglianza che possono riflettere relazioni evolutive, funzionali o strutturali tra gli organismi.

In termini pratici, un allineamento consiste nel disporre le sequenze una sopra l'altra, introducendo eventuali spazi vuoti ("gap") per massimizzare il numero di corrispondenze tra le basi.
Questa operazione consente di identificare mutazioni, inserzioni o delezioni, e di dedurre l'origine comune di geni o proteine.
Queste analisi sono fondamentali in vari ambiti applicativi, tra cui
\begin{itemize}
    \item l'annotazione di nuovi genomi
    \item la ricerca di omologie tra geni
    \item la previsione della funzione di proteine sconosciute
    \item lo studio delle malattie genetiche
\end{itemize}

Esistono due principali categorie di allineamento: globale e locale.
L'allineamento globale, proposto da Needleman e Wunsch (1970), confronta le sequenze nella loro interezza, ed è più adatto quando le sequenze hanno lunghezze simili e un alto grado di somiglianza.
L'allineamento locale, introdotto da Smith e Waterman (1981), mira invece a trovare le regioni più simili all'interno di sequenze potenzialmente molto diverse, risultando particolarmente utile in contesti biologici reali, dove solo alcune porzioni delle sequenze mostrano omologia.

\subsection{Obiettivo del progetto}
L'obiettivo del progetto è implementare e ottimizzare l'algoritmo di Smith-Waterman sfruttando la potenza di calcolo parallelo delle GPU attraverso la piattaforma CUDA.
L'algoritmo di Smith-Waterman presenta una complessità computazionale di tipo quadratico rispetto alla lunghezza delle sequenze da confrontare.
Questo lo rende poco adatto all'elaborazione di grandi quantità di dati, come quelli generati dalle moderne tecniche di sequenziamento del DNA.

L'obiettivo principale è quindi ridurre i tempi di esecuzione attraverso una parallelizzazione efficiente delle sue operazioni e mantenendo inalterata la precisione dell'algoritmo.

Il progetto mira a dimostrare come l'uso del calcolo parallelo possa rendere praticabile l'utilizzo dell'algoritmo di Smith-Waterman anche in contesti reali e su dataset di grandi dimensioni.

\subsection{Strumenti e ambiente di sviluppo}
Lo sviluppo del progetto è stato condotto utilizzando il linguaggio CUDA C, sfruttando l'integrazione diretta con l'architettura di calcolo parallelo delle GPU NVIDIA.

L'ambiente di sviluppo è stato configurato su sistema operativo Windows, utilizzando il compilatore \texttt{nvcc} in versione 12.9.86, incluso nel CUDA Toolkit 13.0.
Il sistema era dotato dei driver NVIDIA in versione 581.57, pienamente compatibili con la versione del toolkit utilizzato.

Per la fase di test e analisi delle prestazioni è stata impiegata una GPU NVIDIA GeForce RTX 3070 (Compute Capability 8.6).
%Per la fase di test e analisi delle prestazioni sono state impiegate due GPU di architettura recente:
%\begin{itemize}
%    \item NVIDIA GeForce RTX 3070 (Compute Capability 8.6)
%    \item NVIDIA GeForce RTX 4060 (Compute Capability 8.9)
%\end{itemize}

\section{Algoritmo}
L'algoritmo di Smith-Waterman permette di ottenere un allineamento locale ottimale tramite un approccio di programmazione dinamica.
In questa sezione viene descritto il funzionamento dell'algoritmo, la sua formulazione e la relativa implementazione sequenziale, che farà da riferimento per la successiva versione parallela su GPU.

\subsection{Algoritmo di Smith-Waterman}
L'algoritmo di Smith-Waterman consente di ottenere un allineamento locale ottimale tra due sequenze, identificando le regioni di massima somiglianza.
Nell'allineamento vengono considerati tre tipi di operazioni: corrispondenze o sostituzioni, inserzioni e delezioni.
Le inserzioni e le delezioni introducono dei "gap" che consentono di allineare correttamente regioni non perfettamente coincidenti tra le sequenze.

L'algoritmo utilizza la tecnica di programmazione dinamica per costruire la matrice di sostituzione in modo da assegnare un punteggio in base ai possibili allineamenti delle sequenze considerate. La soluzione viene poi ricostruita partendo dall'allineamento con il punteggio più alto.

Si articola in quattro fasi principali:
\begin{itemize}
    \item \textbf{Definizione del modello di punteggio e penalità}\\
    Ogni coppia di basi (o amminoacidi, nel caso di sequenze proteiche) riceve un punteggio in base alla corrispondenza: le corrispondenze esatte vengono premiate con valori positivi, mentre le sostituzioni ottengono punteggi inferiori o negativi.
    La scelta dello schema di penalità dipende dal tipo di sequenza analizzata e dagli obiettivi dello studio.

    \item \textbf{Inizializzazione della matrice dei punteggi}\\
    Si costruisce una matrice bidimensionale di dimensioni pari a $(m+1) \times (n+1)$, dove m e n rappresentano le lunghezze delle due sequenze.
    Tutti gli elementi della prima riga e della prima colonna vengono inizializzati a zero.
    Questa scelta consente all'allineamento di iniziare in qualsiasi posizione delle sequenze e garantisce che i gap iniziali non siano penalizzati.

    \item \textbf{Calcolo dei punteggi (fase di riempimento della matrice)}\\
    $$
    H_{i,j} = \max
    \begin{Bmatrix}
    0 \\
    H_{i-1,j-1} + S(a_i, b_j) \\
    H_{i-1,j} + G \\
    H_{i,j-1} + G
    \end{Bmatrix}
    $$
    Per ogni cella si considerano tre possibili punteggi:
    \begin{itemize}
        \item il punteggio diagonale (corrispondente a match/mismatch),
        \item il punteggio proveniente da sinistra (introduzione di un gap nella sequenza verticale),
        \item il punteggio proveniente dall'alto (introduzione di un gap nella sequenza orizzontale).
    \end{itemize}
    Il valore della cella viene impostato come il massimo tra questi tre valori e zero, così da eliminare i contributi negativi e limitare l'allineamento a regioni localmente simili.

    \item \textbf{Traceback (ricostruzione dell'allineamento)}\\
    Una volta completato il riempimento della matrice, si identifica la cella con il punteggio massimo.
    A partire da essa, si effettua un traceback seguendo la direzione dei punteggi che hanno contribuito al valore corrente (diagonale, sinistra o alto), fino a raggiungere una cella con valore zero.
    Il percorso così ottenuto rappresenta l'allineamento locale ottimale.
    È inoltre possibile individuare allineamenti secondari ripetendo il processo di traceback a partire dai punteggi massimi rimanenti.
\end{itemize}

Grazie a questa formulazione l'algoritmo di Smith-Waterman garantisce la massima accuratezza possibile nell'allineamento locale.
Tuttavia, la sua complessità computazionale quadratica in funzione della lunghezza delle sequenze ($O(m \cdot n)$) ne limita l'applicabilità diretta su dataset di grandi dimensioni, rendendo necessarie strategie di ottimizzazione e parallelizzazione.

\subsection{Implementazione sequenziale}
\begin{lstlisting}
void cpu_alignment(
                sw_mat* h,
                string seq_1,
                string seq_2,
                int seq_1_length,
                int seq_2_length
                ) {
    for(int i = 1; i < seq_1_length; ++i) {
        for (int j = 1; j < seq_2_length; ++j) {
            bool match = seq_1[i] == seq_2[j];
            int d = h[INDEX(i - 1, j - 1, seq_2_length)].value +
                            (match ? MATCH : MISMATCH);
            int u = h[INDEX(i - 1, j, seq_2_length)].value + GAP;
            int l = h[INDEX(i, j - 1, seq_2_length)].value + GAP;
            int result = max(max(0, d), max(u, l));
            
            sw_mat* h_current = &h[INDEX(i, j, seq_2_length)];
            h_current->value = result;

            if (result == 0)    h_current->dir = STOP;
            else if (result == d)   h_current->dir = DIAG;
            else if (result == u)   h_current->dir = UP;
            else if (result == l)   h_current->dir = LEFT;
        }
    }
}
\end{lstlisting}

Per questa implementazione sequenziale sono stati scelti i seguenti punteggi e penalità secondo basi empiriche e di pratica comune:
\begin{itemize}
    \item corrispondenza 3
    \item differenza -3
    \item gap -2
\end{itemize}

La matrice è implementata in un array monodimensionale, per garantire coerenza e coesione con i successivi algoritmi paralleli, e l'accesso avviene attraverso una macro, definita di seguito, che linearizza l'indice a partire da riga, colonna e lunghezza della sequenza disposta lungo le colonne.

\begin{lstlisting}
#define INDEX(row, col, stride) ( (row) * (stride) + (col) )
\end{lstlisting}

Le due sequenze sono interscambiabili, in quanto l'obiettivo è quello di trovare degli allineamenti locali, quindi ambivalenti, ma nel codice si è scelto, per motivi di test e chiarezza, di rifersi alle sequenze con nomi quali \texttt{main\_sequence} e \texttt{sequence} (oppure come \texttt{sequence\_1} e \texttt{sequence\_2}), mantenendo sempre la \texttt{main\_sequence} (o \texttt{sequence\_1}, solitamente la stringa più lunga) un carattere per riga e la \texttt{sequence} (o \texttt{sequence\_2}) un carattere per colonna.

La matrice è, quindi, strutturata come array di \texttt{struct sw\_mat} definita come struttura ausiliaria composta da un elemento numerico che rappresenta il punteggio dell'elemento e un enumerazione che rappresenta la direzione da cui deriva il punteggio, cioè \texttt{DIAG, UP, LEFT} come la direzione dell'elemento che ha fornito la combinazione di punteggio massima, oppure \texttt{STOP} ad indicare una casella vuota. La scansione e la popolazione della matrice avviene da sinistra verso destra e dall'alto verso il basso e ogni cella dipende dal valore precedente sulla stessa riga, precedente sulla stessa colonna e precedente sulla stessa diagonale, in modo da valutare la combinazione di punteggio maggiore.

Non verrà discussa in maniera approfondita la logica di traceback, in quanto non è parte dell'argomento di studio. L'algoritmo che ricostruisce la soluzione, come discusso nei capitoli precedenti, parte dal valore maggiore della matrice e ricostruisce la soluzione seguendo il campo \texttt{dir} della struct di cui è composta la matrice fino al primo \texttt{STOP}.

Per valutare la correttezza degli algoritmi paralleli, però, non viene presa in considerazione la soluzione, ma la matrice stessa per le seguenti ragioni: è più semplice e diretto valutare l'uguaglianza di due matrici; non sempre la soluzione è unica; è possibile arrivare alla stessa soluzione apparente (sequenza di basi) seguendo due percorsi diversi.
In questo contesto la correttezza è determinata dagli stessi valori e dallo stesso percorso delle soluzioni, quindi dalla stessa matrice.

\section{Su GPU}
\subsection{Parallelizzazione dell'algoritmo}

L'algoritmo di Smith-Waterman, nella sua formulazione classica (come discusso nel Capitolo 2), si basa sulla programmazione dinamica per calcolare il punteggio di allineamento $H_{i,j}$ per ogni cella della matrice $H$.
La natura della programmazione dinamica introduce una dipendenza rigorosa nel calcolo di $H_{i,j}$: il valore di ogni cella dipende dai valori calcolati nelle celle adiacenti precedentemente completate (in alto, a sinistra e in diagonale).

Questa interdipendenza impedisce il calcolo parallelo dell'intera matrice $H$, rendendo l'algoritmo naturalmente sequenziale nella sua implementazione. Per sfruttare le architetture parallele, è necessario identificare e strutturare la computazione in modo da massimizzare l'indipendenza e il parallelismo.

Esistono due approcci principali per affrontare questa dipendenza e parallelizzare l'algoritmo:
\begin{itemize}
    \item \textbf{Indipendenza antidiagonale:} Osservando la matrice, si nota che tutte le celle situate sulla stessa antidiagonale dipendono unicamente dalle celle calcolate sulle antidiagonali precedenti. Quindi, tutte le celle su una singola antidiagonale possono essere calcolate in parallelo.
    %Questa strategia richiede $O(m+n)$ passaggi seriali, dove $m$ e $n$ sono le lunghezze delle sequenze. In ogni passaggio, $O(\min(i, m+n-i))$ celle possono essere elaborate simultaneamente.
    \item \textbf{Indipendenza della singola colonna:} Questo approccio tenta di parallelizzare il calcolo degli elementi all'interno di una singola colonna (o riga). Sebbene ogni cella $[i,j]$ dipenda intrinsecamente da $[i-1, j]$ (elemento immediatamente sopra) e $[i, j-1]$ (elemento immediatamente a sinistra), l'analisi dell'algoritmo di Smith-Waterman con l'uso di ottimizzazioni può permettere di ristrutturare il calcolo per minimizzare o eliminare la dipendenza verticale tra le righe $i$ e $i-1$ per una colonna fissa $j$. L'obiettivo è calcolare tutti gli elementi della colonna $j$ in parallelo, una volta che la colonna $j-1$ è stata completamente calcolata.
\end{itemize}

Nel contesto di questo progetto di calcolo su GPU, si è scelto di adottare la strategia basata sull'indipendenza della singola colonna. Questa scelta è motivata da diversi fattori pratici e computazionali:

\begin{itemize}
    \item \textbf{Massimizzazione del parallelismo per colonna:} Assegnare un thread (o un blocco di thread) alla computazione di tutti gli elementi di una colonna offre un elevato e uniforme grado di parallelismo, specialmente per sequenze lunghe, semplificando la gestione del carico di lavoro e la mappatura tra thread e dati.
    \item \textbf{Sfruttamento della coalescenza:} L'accesso ai dati lungo una colonna è fondamentale per le prestazioni su GPU. Mantenere l'accesso ai dati in modo sequenziale favorisce la coalescenza degli accessi alla memoria globale.
    \item \textbf{Semplicità di sincronizzazione:} L'intera operazione si riduce a $n$ passaggi sequenziali (il numero di colonne). Tra un passaggio e l'altro è necessaria una barriera di sincronizzazione, rendendo la logica di controllo più semplice rispetto alla complessa gestione delle sincronizzazioni tra richiesta dall'approccio antidiagonale puro.
\end{itemize}

L'approccio della soluzione si basa sull'iterazione seriale sulle colonne della matrice e il calcolo parallelo degli elementi lungo la colonna. L'elemento a sinistra ($H_{i, j-1}$) è già disponibile nella matrice risultante dal lancio del kernel precedente, così come l'elemento sulla diagonale ($H_{i-1, j-1}$), poiché disponibili nella colonna precedente.
La cella superiore, invece, è la dipendenza più critica. Per eliminarla e consentire il calcolo in parallelo lungo la colonna, si fa affidamento sul concetto di gap e sul fatto che il valore stesso della cella superiore sia necessariamente determinato da dei gap verticali a partire da un elemento diagonale o iniziale, in quanto il valore sulla colonna attuale deriva da una combinazione di spostamenti verticali e diagonali.
Il calcolo si riconduce a una catena di dipendenze verso l'alto e prosegue verso l'alto fino a quando si raggiunge un elemento $H_{k,j}$ la cui dipendenza non è verso l'alto, ma verso la cella diagonale (ovvero, c'è un match tra i caratteri), oppure si raggiunge l'inizio della matrice \cite{chen_new_2010}.
Questo permette di ricavare il valore da una cella $H_{k-1, j-1}$, della colonna precedente, purché $S_1[k] = S_2[j]$ sia il primo match trovato salendo da $i$ a $1$:
$$
U = H_{k-1, j-1} + (i-k) \cdot \text{GAP} + \text{MATCH}
$$
Questo calcolo, pur essendo più intensivo in termini di istruzioni, elimina la dipendenza verticale tra $H_{i,j}$ e $H_{i-1, j}$, rendendo tutti i calcoli di lungo la colonna indipendenti gli uni dagli altri.

\subsection{Implementazione su GPU}
\begin{enumerate}
    \item \textbf{Serializzazione sulle Colonne ($j$):} Il loop esterno che avanza colonna per colonna è gestito sul lato CPU (Host). Ad ogni iterazione $j$, un kernel GPU viene lanciato per calcolare l'intera colonna $j$.
    \item \textbf{Parallelizzazione sulle Righe ($i$):} All'interno del kernel, il thread con ID $TID$ viene mappato all'elemento $H_{i,j}$ con $i = TID + 1$. Tutti i thread calcolano simultaneamente i loro rispettivi $H_{i,j}$.
\end{enumerate}
\subsection{Ottimizzazione con blocchi}
\subsection{Ottimizzazione con memoria condivisa (SMEM)}
\subsection{Ottimizzazione con warp e coalescenza}
\subsection{Analisi delle prestazioni e comparazione}
\section{Conclusione}

\newpage
\bibliographystyle{plain}
\bibliography{bibliography}
\end{document}